\subsection{A complete two-class unitary comparison over \texorpdfstring{$\mathbf Q(\sqrt{-7})$}{Q(sqrt(-7))}}
\label{sec:u3-minus7}

\begin{theorem}[K17-U3-MINUS7]\label{thm:u3-minus7-unit}
Put
\[
 E=\mathbf Q(\omega),\qquad \omega^2-\omega+2=0,
 \qquad \omega=\frac{1+\sqrt{-7}}2,
 \qquad q=25078104767246023.
 \tag{U7.1}
\]
Let $G$ be the unitary group of the standard positive Hermitian form
on $E^3$, and let $K=G(\widehat{\mathbf Z})$ be the full integral
lattice level. On the \emph{entire} space $\mathcal A_{15}$ of
algebraic modular forms of highest weight $(15,0,-15)$ and level $K$,
let $T_2^+$ be the positive minuscule Hecke operator, with compact
subgroup volume one, at the embedding $E\hookrightarrow\mathbf Q_2$
where $\omega$ has positive valuation. Then
\[
 \det\bigl((1+2^{16}+2^{32})I-2^{15}T_2^+
                  \mid\mathcal A_{15}\bigr)
 \quad\hbox{is a unit at every place above }q.
 \tag{U7.2}
\]
The same assertion holds for the opposite minuscule direction.
The eigenvalues of $2^{15}T_2^+$ are $2^{16}$ times the standard
normalized Satake trace. This is a finite Hecke-model theorem;
coverage of an arithmetic source family is a separate assertion.
\end{theorem}

\begin{proof}
We give the complete class-set input, actual coset construction,
integral specialization argument, and finite certificate. The last
step works on a $171$-dimensional space and does not require a basis
of the invariants of the second arithmetic group.

\paragraph{The two adelic classes and their stabilizers.}
For the group and level in the statement, the complete class set has
representatives $1,\alpha$, where $\alpha$ is the identity away $2$
and, at the chosen embedding, is represented by
\[
 \alpha_2=\begin{pmatrix}1&0&0\\0&1&0\\-1&-1&2\end{pmatrix}.
 \tag{U7.3}
\]
The stabilizers
$\Gamma_1=G(\mathbf Q)\cap K$ and
$\Gamma_2=G(\mathbf Q)\cap\alpha K\alpha^{-1}$ have orders
$48$ and $336$. Here $\Gamma_1$ is the full group of signed
permutation matrices and $\Gamma_1\subset\Gamma_2$ has index $7$.
More explicitly, $\Gamma_2$ is generated by $\Gamma_1$ and
\[
 \tau=\frac12
 \begin{pmatrix}-1&1-\omega&1\\-1&-1+\omega&1\\
                 \omega&0&\omega\end{pmatrix}.
 \tag{U7.4}
\]
These are the exact class-set calculations in
\cite[\S3, printed pp.~6--7]{K17Loeffler2008}. Their completeness
uses the mass $1/42=1/48+1/336$, not a restriction to principal
lattices. Evaluation at the two representatives identifies the full
space of algebraic modular forms with
\[
 W_1\oplus W_2,
 \qquad W_i=\mathcal H_{15}^{\Gamma_i},
 \qquad W_2\subset W_1.
 \tag{U7.5}
\]
Here $\mathcal H_k$ is the harmonic model described below. We use
the convention $f(\gamma gk)=\rho(\gamma)f(g)$, as in
\cite[\S\S2.1--2.2, printed pp.~3--4]{K17Loeffler2008}.

\paragraph{The complete actual correspondence at $2$.}
The double coset of local type $(0,0,1)$ has $2^2+2+1=7$ right
cosets. On the class $1$, six are principal and the seventh is
$\alpha K$ itself; on the class $\alpha$, all seven are principal.
These assertions are established in
\cite[Lemma~2.5 and \S3, printed pp.~5--7]{K17Loeffler2008}.
For full reproducibility we recover the rational representatives by
the following exact finite enumeration, which also checks every
elementary-divisor condition.

Write $\overline\omega=1-\omega$. For $r=2,4$, enumerate all
unordered orthogonal triples of columns in $\mathbf Z[\omega]^3$
of squared norm $r$, identifying a column with its negative. If a
scalar is $a+b\omega$, its norm is
\[
 a^2+ab+2b^2=(a+b/2)^2+7b^2/4.
 \tag{U7.6}
\]
Consequently $|b|\le1$, $|a|\le3$ exhaust every scalar of norm at
most $4$. A canonical sign and lexicographic ordering eliminate
exactly the right action of $\Gamma_1$. If $d=a+b\omega$ is the
determinant of the column matrix, retain precisely the frames with
the following properties:
\[
 \begin{array}{c|l|r|r|r}
 r&\text{retention condition}&\text{columns}
                &\text{all frames}&\text{retained}\\ \hline
 2&a+2b\equiv2\pmod4&12&14&6\\
 4&a+6b\equiv4\pmod8,
       \;\sum_i a_{ij}\equiv0\pmod2\ \text{for every }j
       &57&113&7.
 \end{array}
 \tag{U7.7}
\]
Here the column count is taken modulo signs, and $a_{ij}$ is the
constant coefficient of the $(i,j)$ entry.
Denote the two resulting sets by $\mathcal U_2,\mathcal U_4$.
Their rational unitary representatives are respectively
\[
 g_U=\frac{U}{1-\omega}\quad(U\in\mathcal U_2),
 \qquad
 g_V=\frac{V}{(1-\omega)^2}\quad(V\in\mathcal U_4).
 \tag{U7.8}
\]
All the equalities $U^*U=2I$, $V^*V=4I$ are checked exactly.

At the chosen $2$-adic embedding, $\omega\equiv2\pmod4$ and
$\omega\equiv6\pmod8$, whereas $1-\omega$ is a unit. Thus the
first determinant condition in (U7.7) says exactly that $g_U$ is
integral with determinant valuation $1$, hence has type $(0,0,1)$.
For the second row, the last row of $\alpha_2^{-1}$ is
$(1/2,1/2,1/2)$. The parity condition says exactly that
$\alpha_2^{-1}g_V$ is integral; its determinant has valuation
$2-1=1$, so it again has type $(0,0,1)$. Unitary duality supplies
the corresponding condition at the conjugate prime. At every
other prime the matrices and their inverses are integral.
Distinct canonical frames give distinct right cosets, since two
rational representatives differ on the right by $K$ exactly when
their ratio is in $\Gamma_1$. There are therefore exactly six
principal cosets in the first row and seven in the second, matching
the full correspondence above.

Set
\[
 S_{11}=\sum_{U\in\mathcal U_2}\rho(g_U),
 \qquad S_{21}=\sum_{V\in\mathcal U_4}\rho(g_V).
\]
The complete Hecke operator is
\[
 T_2^+=\begin{pmatrix}S_{11}&\operatorname{inc}\\
                           S_{21}&0\end{pmatrix}
       \quad\text{on }W_1\oplus W_2.
 \tag{U7.9}
\]
In particular $S_{21}$ maps $W_1$ into $W_2$: this follows from the
actual Hecke correspondence and its equivariance. It is not an
assumption inferred from finitely many evaluations. The upper-right
entry is exactly inclusion because the seventh representative in
the first row is $\alpha$ with rational multiplier $1$.

\paragraph{The coefficient lattice and its reduction.}
Take independent variables $z_1,z_2,z_3,w_1,w_2,w_3$ and put
\[
 R=\sum_i z_iw_i,\qquad
 C=\sum_i\partial_{z_i}\partial_{w_i},\qquad
 \rho(g)f(z,w)=f(g^{-1}z,\overline g^{-1}w).
 \tag{U7.10}
\]
The kernel of $C$ in bidegree $(k,k)$ is the full irreducible
representation of highest weight $(k,0,-k)$. One can verify this
without a decomposition calculation: $C$ is surjective, its kernel
has dimension
$\binom{k+2}{2}^2-\binom{k+1}{2}^2=(k+1)^3$, and it contains
the highest-weight vector $w_1^kz_3^k$, whose Weyl dimension is
$(k+1)^3$.

The following explicit right inverse also controls specialization.
For $F$ of bidegree $(k-1,k-1)$,
\[
 C\bigl(R^{j+1}C^jF\bigr)
 =(j+1)(2k+1-j)R^jC^jF+R^{j+1}C^{j+1}F.
\]
Thus $\sum_{j=0}^{k-1}c_jR^{j+1}C^jF$ is a preimage of $F$ when
\[
 c_0=\frac1{2k+1},\qquad
 c_j=-\frac{c_{j-1}}{(j+1)(2k+1-j)}.
 \tag{U7.11}
\]
For $k=15$ every prime in these denominators is at most $31$.
The certificate proves that $q$ is prime by the full-order Lucas
witness $3$, using
\[
 q-1=2\cdot3^3\cdot10487\cdot44284289839,
 \tag{U7.12}
\]
with primality of the factors proved by trial division. It also
checks $(-7)^{(q-1)/2}\equiv-1\pmod q$. Therefore
\[
 A=\mathbf Z_{(q)}[\omega]
 \quad\text{has residue field}\quad
 \mathbf F_{q^2}=\mathbf F_q[\omega]/(\omega^2-\omega+2).
 \tag{U7.13}
\]
The contraction is split surjective over $A$ by (U7.11). Its
harmonic kernel is a free direct summand and reduction of the
kernel is the kernel after reduction. All matrices in $\Gamma_2$
have denominators which are powers of $2$, by (U7.4), so both
arithmetic groups preserve this coefficient lattice. Their orders
$48,336$ are units in $A$. Averaging therefore shows that taking
their invariants commutes with reduction. Write
$M_i=\mathcal H_{15,A}^{\Gamma_i}$. These are free $A$-modules,
$M_2\subset M_1$, and (U7.9) defines an integral operator on
$M_1\oplus M_2$. These facts exclude any change of harmonic or
invariant dimension at $q$.

\paragraph{The entire $171$-dimensional first invariant space.}
In bidegree $(k,k)$ the $\Gamma_1$-invariant polynomials have the
following complete orbit-sum basis:
\[
 |a|=|b|=k,\qquad a_i-b_i\equiv0\pmod2,
 \qquad
 \sum_{\text{distinct simultaneous permutations}}z^a w^b.
 \tag{U7.14}
\]
The parity condition imposes invariance under each coordinate sign;
the orbit sum imposes the full permutation group. These basis
vectors have disjoint monomial supports. Exhausting all pairs of
three-part compositions gives $795$ source orbits at $k=15$ and
$624$ target orbits at $k=14$. The contraction is computed from
\[
 C(z^aw^b)=\sum_i a_i b_i z^{a-e_i}w^{b-e_i}.
\]
The saved $171$ columns are annihilated by this full matrix modulo
$q$, and have $171$ distinct free coordinates forming the identity
matrix. They are therefore independent. Surjectivity of (U7.11),
including after taking invariants, proves that they span the entire
kernel of dimension $795-624=171$. The certificate does not assume
that a selected collection of forms exhausts the coefficient space.

\paragraph{Schur elimination and exact denominator cancellation.}
Let $k=15$ and put
\[
 E_0=1+2^{16}+2^{32}=4295032833.
\]
This is a unit in $A$. From (U7.9), an element in the kernel of
$E_0I-2^kT_2^+$ satisfies
$v_2=E_0^{-1}2^kS_{21}v_1$. Consequently the full operator is
invertible if and only if the following endomorphism of $M_1$ is
invertible:
\[
 D=E_0^2I-E_0\,2^kS_{11}
                  -2^{2k}\operatorname{inc}S_{21}.
 \tag{U7.15}
\]
This is ordinary block elimination over $A$ and over its residue
field; all scalar factors removed are units. No knowledge of
$\dim M_2$ is required.

For $U^*U=2I$ and $V^*V=4I$, homogeneity in the two sets of
variables gives
\[
 \begin{split}
 2^k\rho(g_U)f(z,w)&=f(U^*z,U^tw),\\
 2^{2k}\rho(g_V)f(z,w)&=f(V^*z,V^tw).
 \end{split}
 \tag{U7.16}
\]
For example $g_U^{-1}=U^*/\omega$ and
$\overline g_U^{-1}=U^t/(1-\omega)$, so the product of the two
homogeneous denominators is $2^k$. The second identity has
denominator $4^k$. Thus every entry needed in (U7.15) is obtained
by integral polynomial substitutions before reduction.

\paragraph{The complete finite certificate.}
Choose the $173$ saved pairs of rational-coordinate points in
$\mathbf F_q^3\times\mathbf F_q^3$, and evaluate all $171$
harmonic basis polynomials. Let $A_0$ be the first $171$ rows of
this evaluation matrix; let $B_2,B_4$ be the corresponding sums
of substituted evaluations in (U7.16). The exact results are
\[
 \begin{array}{c|r}
 \text{quantity in }\mathbf F_{q^2}&\text{value in }\mathbf F_q\\ \hline
 \det A_0&9611653439256081\\
 \det(E_0^2A_0-E_0B_2-B_4)&18272413515919144\\
 \det D&13187546903051047.
 \end{array}
 \tag{U7.17}
\]
The first determinant is nonzero, so evaluation is injective on
the entire reduced $M_1$. Since the actual correspondence already
preserves $M_1$, its endomorphism is identified by these values.
The last determinant is the quotient of the first two appropriate
determinants and is also nonzero. This proves invertibility of
$D$ modulo $q$, hence (U7.2) by (U7.15) and the lattice argument.
The two additional point rows verify the same recovered operator
independently of the identifying rows. These extra checks are not
used as a substitute for injectivity or Hecke equivariance.

Here is the complete replay specification; all files are in
\path{work/hecke_minus7/}.
\begin{enumerate}
\item \path{neighbor_frames.py} exhausts the bounded integral
      enumeration in (U7.6)--(U7.8), saving every retained frame
      in \path{neighbor_frames.json}.
\item \path{prepare_model.py} constructs all orbit sums, the full
      contraction, its kernel, and the exact points, saving
      \path{model_k15.json} and the matching integer text input.
\item \path{compute_schur.cpp} uses only exact residue arithmetic
      and writes all evaluation arrays and the complete recovered
      $171$ by $171$ operator to \path{schur_k15.json}.
\item The independent standard-library verifier
      \path{verify_schur.py} reconstructs the prime proof, field,
      every frame class, all source and target orbits, and the full
      harmonic kernel without importing the preparation code. It
      proves equality of the integer text and JSON inputs,
      recomputes every polynomial evaluation at all $173$ points
      for all $13$ frames, checks the entire operator identity,
      and independently computes all three determinants by
      column elimination. Its exact replay reports PASS in
      \path{verify_schur.json}, with input hashes and elapsed time.
\end{enumerate}
The C++ replay takes approximately $1.1$ seconds and the independent
Python replay approximately $18.6$ seconds on the present executor.
All field operations use integers; time measurements play no role
in the certificate.

\paragraph{Direction and Satake normalization.}
Complex conjugation on $G(\mathbf A_f)$ preserves $K$ and exchanges
the two minuscule directions. The coefficient involution
$Jf(z,w)=f(w,z)$ satisfies
$J\rho(\overline\gamma)=\rho(\gamma)J$.
Hence $f(g)\mapsto J f(\overline g)$ is an automorphism of the
full algebraic modular-form space conjugating the two Hecke
operators. It is integral at $q$: both class representatives are
trivial away $2$, so their conjugate-class comparison multipliers
are integral at $q$. The assertion therefore holds for either
direction, without assuming that each individual form is fixed
by complex conjugation.

For $\mathrm{GL}_3$, the positive minuscule characteristic function
has normalized Satake transform $p$ times the standard trace;
see \cite[\S3, (3.13)--(3.14)]{K17GrossSatake1998}.
Thus $2^{15}T_2^+$ has eigenvalues $2^{16}$ times that trace,
as asserted.
\end{proof}
